\section{Introducción}

En las últimas décadas, los sistemas inteligentes y los agentes autónomos han experimentado un gran crecimiento dentro del ámbito de la Inteligencia Artificial. Estos agentes son capaces de percibir su entorno, razonar sobre la información disponible y ejecutar acciones orientadas al cumplimiento de determinados objetivos. Tradicionalmente, gran parte de los modelos desarrollados en sistemas multiagente se han basado en enfoques puramente racionales, donde la toma de decisiones depende principalmente de procesos lógicos y de optimización.

Sin embargo, numerosos estudios en psicología, neurociencia y ciencias cognitivas han demostrado que las emociones desempeñan un papel fundamental en la cognición humana, influyendo directamente en la memoria, el razonamiento, la percepción y la toma de decisiones. Lejos de representar únicamente respuestas irracionales, las emociones permiten priorizar información, adaptar el comportamiento al contexto social y mejorar los procesos deliberativos en entornos complejos e inciertos \cite{alfonso2021affect}.

Como consecuencia, ha surgido un creciente interés en incorporar mecanismos emocionales dentro de agentes inteligentes y sistemas multiagente. Este campo, conocido como \textit{computación afectiva}, busca desarrollar sistemas capaces de representar, interpretar y utilizar emociones durante sus procesos internos de razonamiento e interacción. El objetivo no es únicamente simular emociones humanas de forma superficial, sino conseguir que los agentes puedan adaptar su comportamiento y reaccionar de manera más coherente en contextos dinámicos o sociales.

Dentro de este contexto, diferentes modelos psicológicos han servido como base teórica para el desarrollo de arquitecturas afectivas computacionales. Entre los modelos más relevantes destacan el modelo OCC (\textit{Ortony, Clore and Collins}), centrado en la evaluación cognitiva de eventos y acciones, cuya formalización lógica ha sido clave para evitar ambigüedades en su implementación computacional \cite{occ2009revisited}, y los modelos dimensionales como PAD (\textit{Pleasure-Arousal-Dominance}), utilizados para representar estados emocionales mediante espacios continuos.

A partir de estas bases teóricas, se han desarrollado múltiples arquitecturas emocionales orientadas a integrar emociones en agentes inteligentes. Entre ellas destacan propuestas como EBDI, una extensión emocional de la arquitectura BDI tradicional \cite{jiang2007ebdi}, así como modelos más recientes centrados en la simulación de emociones primarias y secundarias \cite{becker2009wasabi}, y en la representación cultural de emociones mediante lógica difusa \cite{taverner2021fuzzy,taverner2023multidimensional}. Asimismo, investigaciones recientes han explorado la relación entre emociones y normas sociales dentro de sistemas multiagente, permitiendo modelar fenómenos como culpa, vergüenza, aprobación social o cooperación normativa \cite{argente2022normative}.

Más allá de presentar modelos aislados, este trabajo analiza cómo las distintas propuestas estudiadas representan una evolución progresiva dentro de la computación afectiva. En particular, el modelo OCC proporciona una estructura cognitiva para la generación de emociones; arquitecturas como EBDI integran dichas emociones dentro del razonamiento deliberativo de agentes BDI; WASABI traslada estos mecanismos a humanos virtuales con interacción expresiva; mientras que los modelos dimensionales y difusos modernos abordan problemas relacionados con la intensidad emocional, la continuidad de estados afectivos y la adaptación cultural.

El presente trabajo tiene como objetivo principal realizar una revisión de los modelos computacionales y arquitecturas utilizadas para la incorporación de emociones en agentes inteligentes. Para ello, el documento se ha estructurado de la siguiente manera: en primer lugar, se abordan los fundamentos teóricos de las emociones y su relación con la computación afectiva. A continuación, se analizan las principales arquitecturas afectivas de la literatura (como EBDI) y los modelos de representación emocional, prestando especial atención a la adaptación cultural. Finalmente, se explora la integración entre emociones y normas sociales en sistemas multiagente, concluyendo con una discusión sobre los retos actuales y las posibles líneas futuras de investigación.